\section{Stable functoriality}

The construction-independent repository source is
\href{../../verified/STABLE_NORMALIZATION_FUNCTORIALITY.md}
{Stable normalization functoriality}.  It also records the relative
differential, Fitting-ideal, fiber-product, and conductor base-change clauses.
The proposition below is the paper-local specialization needed for the
boundary signature.

\begin{proposition}[stable normalization functoriality]
\label{prop:stable-normalization}
Let $Y$ be a normal integral affine variety over $\C$, let $L/k(Y)$ be a
finite extension, and put $X=\Norm_Y(L)$.  Let $U\subset X$ be a dense open
and $\partial=(X\setminus U)_{\mathrm{red}}$.  For $s\ge0$, write
\[
 Y_s=Y\times\A^s,\qquad U_s=U\times\A^s,
 \qquad L_s=L(t_1,\ldots,t_s).
\]
Then there is a canonical isomorphism over $Y_s$
\[
 \Norm_{Y\times\A^s}\bigl(L(t_1,\ldots,t_s)\bigr)
 \simeq X\times\A^s.                                    \tag{5.1}
\]
In particular, taking $L=k(U)$ and $X=\overline X_F$ gives
\[
 \Norm_{Y\times\A^s}\bigl(k(U)(t_1,\ldots,t_s)\bigr)
 \simeq\overline X_F\times\A^s.
\]
Under (5.1) the following statements hold.
\begin{enumerate}
\item The reduced boundary is
\[
 (X\times\A^s\setminus U_s)_{\mathrm{red}}
 =\partial\times\A^s,                                   \tag{5.2}
\]
and its prime divisors are exactly $E\times\A^s$, with $E$ a
prime divisor of $\partial$.
\item If a boundary prime $E$ lies over a target prime divisor $Z$, the
corresponding stabilized primes are $E_s=E\times\A^s$ and
$Z_s=Z\times\A^s$.  Their divisorial valuations are the Gauss extensions
\[
 v_{E_s}|_L=v_E,\qquad v_{E_s}(t_i)=0,
\]
and
\[
 e(E_s/Z_s)=e(E/Z),\qquad f(E_s/Z_s)=f(E/Z).             \tag{5.3}
\]
\item Every reduced or scheme-theoretic multiple intersection commutes with
stabilization.  In particular, for target divisors $Z_1,\ldots,Z_q$ and
boundary divisors $E_1,\ldots,E_p$,
\[
 \bigcap_i(Z_i\times\A^s)
 =\left(\bigcap_i Z_i\right)\times\A^s,
 \qquad
 \bigcap_j(E_j\times\A^s)
 =\left(\bigcap_j E_j\right)\times\A^s                 \tag{5.4}
\]
as schemes.
\item If $R$ is the affine coordinate ring of one of these intersections,
then the stabilized ring is $R[t_1,\ldots,t_s]$ and
\[
 \Nil(R[t_1,\ldots,t_s])=\Nil(R)[t_1,\ldots,t_s].       \tag{5.5}
\]
Consequently reducedness and the exact nilpotency index of the nilradical
are preserved.
\end{enumerate}
\end{proposition}

\begin{proof}
This is the specialization of the construction-independent
\href{../../verified/STABLE_NORMALIZATION_FUNCTORIALITY.md}
{stable-normalization functoriality theorem}.  Apply its normalization,
boundary, valuation, intersection, and nilradical clauses to the pair
$(X,U)$.
\end{proof}

\begin{lemma}[stable boundary signature]\label{lem:stable}
For the maps in Lemma~\ref{lem:boundary}, the pair $(e_\Delta,\mu)$ of
(2.1) is invariant under polynomial left--right equivalence and after
adjoining any number of identity variables.
\end{lemma}

\begin{proof}
A left--right equivalence identifies the corresponding finite function-field
extensions over isomorphic targets.  Uniqueness and functoriality of
normalization identify the finite covers, carry the distinguished affine
open to the distinguished affine open, and preserve boundary valuations,
their $e,f$ labels, and scheme-theoretic target intersections.  The intrinsic
ramification marking from Lemma~\ref{lem:boundary} prevents the two target
vertices from being exchanged.

For stabilization, Proposition~\ref{prop:stable-normalization} identifies the
normalization and the distinguished open pair, preserves all reduced
boundary primes and their valuation labels, and base-changes the marked
target intersection scheme.  Its final assertion preserves reducedness and
the exact nilpotency index.  Thus both entries of $(e_\Delta,\mu)$ are
unchanged.
\end{proof}
